% Part: first-order-logic
% Chapter: syntax-and-semantics
% Section: Structures

\documentclass[../../../include/open-logic-section]{subfiles}

\begin{document}

\olfileid[pt]{fol}{syn}{str}

\olsection{\printtoken{P}{structure} para Linguagens de Primeira Ordem}

\begin{explain}
As linguagens de primeira ordem são, por si mesmas, \emph{não interpretadas:} os
!!{constant}s, !!{function}s e !!{predicate}s não têm nenhum significado
específico associado a eles. Os significados são dados especificando-se
\article{structure} \emph{!!{structure}}. Ela especifica o
\emph{domínio}, isto é, os objetos que os !!{constant}s denotam, sobre os quais os
!!{function}s operam e sobre os quais os quantificadores variam. Além disso,
ela especifica quais !!{constant}s denotam quais objetos, como
!!a{function} mapeia objetos para objetos, e a quais objetos os
!!{predicate}s se aplicam. !!^{structure}s são a base para as
noções \emph{semânticas} em lógica, por exemplo, a noção de consequência,
validade, satisfatibilidade. Elas são chamadas de várias formas na
literatura: ``estruturas'', ``interpretações'' ou ``modelos''.
\end{explain}

\begin{defn}[!!^{structure}s]
\Article{structure} \emph{!!{structure}}~$\Struct M$, para uma linguagem
$\Lang{L}$ de lógica de primeira ordem consiste nos seguintes elementos:
\begin{enumerate}
\item \emph{Domínio:} um conjunto não vazio, $\Domain M$
\item \emph{Interpretação dos !!{constant}s:} para cada !!{constant}~$c$ de
  $\Lang{L}$, !!a{element} $\Assign{c}{M} \in \Domain M$
\item \emph{Interpretação dos !!{predicate}s:} para cada
  !!{predicate}~$R$ de $n$ lugares de $\Lang{L}$ (distinto de $\eq$), uma relação
  de $n$ lugares $\Assign{R}{M} \subseteq \Domain{M}^n$
\item \emph{Interpretação dos !!{function}s:} para cada
  !!{function}~$f$ de $n$ lugares de $\Lang{L}$, uma função de $n$ lugares $\Assign{f}{M}
  \colon \Domain{M}^n \to \Domain{M}$
\end{enumerate}
\end{defn}

\begin{ex}
!!^a{structure}~$\Struct M$ para a linguagem da aritmética consiste em um
  conjunto, um elemento de $\Domain M$, $\Assign{\Obj 0}{M}$, como
  interpretação do !!{constant}~$\Obj 0$, uma função de um lugar
  $\Assign{\Obj \prime}{M} \colon \Domain{M} \to \Domain M$, duas
  funções de dois lugares $\Assign{\Obj +}{M}$ e $\Assign{\Obj
    \times}{M}$, ambas $\Domain M^2 \to \Domain M$, e uma relação de
  dois lugares $\Assign{\Obj <}{M} \subseteq \Domain{M}^2$.

Um exemplo óbvio de tal estrutura é o seguinte:
\begin{enumerate}
\item $\Domain N = \Nat$
\item $\Assign{\Obj 0}{N} = 0$
\item $\Assign{\Obj \prime}{N}(n) = n + 1$ para todos $n \in \Nat$
\item $\Assign{\Obj +}{N}(n, m) = n + m$ para todos $n, m \in \Nat$
\item $\Assign{\Obj \times}{N}(n, m) = n\cdot m$ para todos $n, m \in \Nat$
\item $\Assign{\Obj <}{N} = \Setabs{\tuple{n, m}}{n \in \Nat, m \in
  \Nat, n < m}$
\end{enumerate}
A estrutura~$\Struct N$ para $\Lang L_A$ assim definida é chamada de
\emph{modelo padrão da aritmética}, porque interpreta as
constantes não lógicas de~$\Lang L_A$ exatamente como se esperaria.

No entanto, há muitas outras !!{structure}s possíveis para~$\Lang
L_A$. Por exemplo, poderíamos tomar como domínio o conjunto~$\Int$ dos
inteiros em vez de~$\Nat$, e definir as interpretações de $\Obj
0$, $\Obj \prime$, $\Obj +$, $\Obj \times$, $\Obj <$ adequadamente.  Mas
também podemos definir estruturas para~$\Lang L_A$ que não têm nada, nem mesmo
remotamente, a ver com números.
\end{ex}

\begin{ex}
Uma estrutura~$\Struct M$ para a linguagem~$\Lang L_Z$ da teoria dos conjuntos requer
apenas um conjunto e uma única relação de dois lugares. Então, tecnicamente, por exemplo, o
conjunto das pessoas mais a relação ``$x$ é mais velho que $y$'' poderia ser
usado como !!a{structure} para $\Lang L_Z$, assim como $\Nat$ junto
com $n \ge m$ para $n, m \in \Nat$.

Uma !!{structure} particularmente interessante para $\Lang L_Z$ na qual os
!!{element}s do domínio são de fato conjuntos, e a interpretação
de $\Obj \in$ é de fato a relação ``$x$ é !!a{element} de~$y$''
é a !!{structure}~$\Struct{{HF}}$ dos \emph{conjuntos hereditariamente finitos}:
\begin{enumerate}
\item $\Domain{{{HF}}} = \emptyset \cup \Pow{\emptyset} \cup
  \Pow{\Pow{\emptyset}} \cup \Pow{\Pow{\Pow{\emptyset}}} \cup \dots$;
\item $\Assign{\Obj \in}{{{HF}}} = \Setabs{\tuple{x, y}}{x, y \in
  \Domain{{{HF}}}, x \in y}$.
\end{enumerate}
\end{ex}

\begin{digress}
As estipulações que fazemos sobre o que conta como !!a{structure} impactam
nossa lógica. Por exemplo, a escolha de impedir domínios vazios garante,
dada a explicação usual de satisfação (ou verdade) para sentenças
quantificadas, que $\lexists[x][(!A(x) \lor \lnot !A(x))]$ é
válida---isto é, uma verdade lógica. E a estipulação de que todos os
!!{constant}s devem referir-se a um objeto no domínio garante que a
generalização existencial é um padrão de inferência correto: $!A(a)$,
portanto $\lexists[x][!A(x)]$. Se permitíssemos que nomes se referissem a algo fora
do domínio, ou que não se referissem, então estaríamos a caminho de uma
\emph{lógica livre}, na qual a generalização existencial requer uma
premissa adicional: $!A(a)$ e $\lexists[x][\eq[x][a]]$, portanto
$\lexists[x][!A(x)]$.
\end{digress}

\end{document}
